% Part: lambda-calculus
% Chapter: syntax
% Section: term-revisited

\documentclass[../../../include/open-logic-section]{subfiles}

\begin{document}

\olfileid[fr]{lam}{syn}{tr}
\olsection{Les termes comme classes d'\texorpdfstring{$\alpha$}{α}-équivalence}

Désormais, nous considérerons les termes à
$\alpha$-équivalence près. Quand nous écrirons
un terme, nous entendrons donc sa classe
d'$\alpha$-équivalence. Par exemple,
$\lambd[a][\lambd[b][a c]]$ désignera
l'ensemble des termes qui lui sont
$\alpha$-équivalents, parmi lesquels
$\lambd[a][\lambd[b][a c]]$,
$\lambd[b][\lambd[a][b c]]$, etc.

Dans les sections précédentes, des lettres
comme $N, Q$ désignaient des termes; elles
désigneront désormais des classes de termes.
Ce sont ces classes, et non leurs représentants
pris isolément, que nous étudierons dans la
suite. Les lettres telles que $x, y$
continueront de désigner des variables.

Nous noterons $\rep{M}$ un \emph{élément}
arbitraire de la classe~$M$, et
$\rep{M}[0]$, $\rep{M}[1]$, etc.,
si nous devons en choisir plusieurs.

Nous réutiliserons les notations des termes
pour alléger les énoncés. Sur les classes,
nous posons les définitions suivantes.
\begin{defn}
\begin{enumerate}
\item $\lambd[x][N]$ est la classe contenant
  $\lambd[x][\rep{N}]$.
\item $PQ$ est la classe contenant
  $\rep{P}\rep{Q}$.
\end{enumerate}
\end{defn}

Ces opérations ne dépendent pas du choix
des représentants, car
l'$\alpha$-conversion est compatible avec
l'abstraction et l'application.

\begin{defn} \ollabel{def:fv}
Les \emph{variables libres} d'une classe
d'$\alpha$-équivalence~$M$, ou $FV(M)$,
sont les variables libres de $\rep{M}$,
c'est-à-dire $FV(\rep{M})$.
\end{defn}

Cette définition ne dépend pas du représentant,
puisque $FV(\rep{M}[0]) =
FV(\rep{M}[1])$, d'après le
\olref[alp]{thm:fv}.

Nous réutilisons aussi la notation de
substitution pour les classes.
\begin{defn} \ollabel{def:sub}
La \emph{substitution} de $R$ à $y$ dans $M$,
notée $\Subst{M}{R}{y}$, est la classe
contenant $\Subst{{\rep{M}}}{\rep{R}}{y}$,
pour tout choix de $\rep{M}$ et de
$\rep{R}$ qui rend cette substitution définie.
\end{defn}

Le \olref[alp]{cor:sub} garantit l'existence de représentants
pour lesquels la substitution est définie. Tous les couples
de représentants qui rendent cette substitution définie
donnent des termes dans une même classe d'$\alpha$-équivalence.

Cette définition simplifie sensiblement les
calculs. Prenons une variable $z$ distincte
de $x$ et de $y$. Alors :
\begin{align}
  \Subst{\lambd[x][x]}{y}{x}
  &\ollabel{eq:1}\\
  &= \Subst{\lambd[z][z]}{y}{x}
     \ollabel{eq:2}\\
  &= \lambd[z][\Subst{z}{y}{x}]\\
  &= \lambd[z][z].
\end{align}

Si l'on considère encore des termes ordinaires,
la substitution de \olref{eq:1} est indéfinie
selon les règles précédentes. Ici, elle porte
sur des classes, ce qui permet le passage
de \olref{eq:2} : nous pouvons remplacer
$\lambd[x][x]$ par $\lambd[z][z]$, puisque
ces deux termes appartiennent à la même classe.

Pour la même raison, nous supposerons désormais
que les représentants choisis satisfont toutes
les conditions nécessaires à la substitution.
Par exemple, devant
$\Subst{\lambd[x][N]}{R}{y}$, nous choisirons
le représentant $\lambd[x][N]$ de sorte que
$x \neq y$, $x \notin FV(R)$ et que la
substitution dans le corps soit définie.

Comme il serait étrange d'appeler
$\lambd[x][x]$ une « classe », nous
appellerons désormais ces objets des
$\Lambda$-termes, ou simplement des « termes »
dans la suite de cette partie, afin de les
distinguer des $\lambda$-termes ordinaires.

Pour définir une opération sur les $\Lambda$-termes à partir
d'une opération sur les représentants, on vérifie deux points.
Il faut pouvoir choisir des représentants pour lesquels le
calcul est défini. Deux choix qui rendent le calcul possible
doivent donner des résultats $\alpha$-équivalents.
La substitution donne un exemple de cette méthode.

\end{document}
